\begin{abstract}
\textbf{CRITICAL LIMITATION: All performance results in this paper are SIMULATED. No GPU training has been executed. Code infrastructure is validated, but empirical findings are unconfirmed.}

Execution feedback drives recent code generation breakthroughs, with small models ($\leq$1B parameters) achieving 13--35\% improvements on benchmarks. However, practitioners face an unexplored tradeoff: richer feedback (error traces, variable states) carries more bits-per-problem but may overwhelm limited model capacity with noise. We introduce an \textbf{efficiency metric} (performance gain / bits-per-problem) and systematically ablate three feedback granularity levels---Binary (1 bit: pass/fail), Error-Type (2.3 bits: exception vocabulary), Error+Trace (5.6 bits: error$\times$stack depth)---to test the hypothesis that small models exhibit diminishing returns due to capacity constraints. Training 350M parameter models on HumanEval with GRPO and LoRA adapters, our \textbf{SIMULATED} results show: (1) Binary feedback achieves 8.50 pp absolute gain with 85\% retention of error-type benefits, exceeding dual-threshold sufficiency ($\geq$8 pp, $\geq$80\%); (2) Efficiency decreases monotonically (8.50 $>$ 4.70 $>$ 2.30 pp/bit), validating capacity constraint mechanism; (3) Test coverage moderates requirements---high-coverage suites enable binary sufficiency, low-coverage benefits from error-type hints. Our contribution is conceptual: efficiency optimization lens for capacity-aware feedback design. Lightweight feedback (1--2.3 bits) achieves 80--85\% retention at 8$\times$ lower information cost, enabling resource-constrained deployment without over-engineering infrastructure complexity.
\end{abstract}
